\section{Analysis of examples}\label{sec: ex} 
In this section, we verify the various claims made in \cref{ex: ternary}. Let $a,b,r,s, F$ and $A, \chi,\gamma,\zeta$ be as in the setup of \cref{ex: ternary}.
\subsection{Failure of the inequality} Notice that $F$ is symmetric in $y,z$. Let $\tilde{p}=\frac{p}{750}=\chi^2-\chi(\gamma+\zeta)-\frac{14}{375}(\gamma^2+\zeta^2)+\frac{49}{50}\gamma\zeta$. When $\epsilon=0$, we have
\begin{align*}
    \deg(\tilde p \chi^2)=&-\frac{7}{250}\\
    \deg(\tilde p \chi\gamma)=\deg(\tilde p\chi\zeta)=&\frac{7}{500}\\
    \deg(\tilde p \gamma^2)=\deg(\tilde p \zeta^2)=&-\frac{17}{300}\\
    \deg(\tilde p \gamma\zeta)=&
\frac{16}{375}\end{align*}
and therefore
\[
    \deg(\tilde p^2)=(-\frac{7}{250})-2\cdot \frac{7}{500}-2\cdot\frac{14}{375}\cdot(-\frac{17}{300})+\frac{49}{50}\cdot \frac{16}{375}=-\frac{56}{5625}
\]
and
\[
    \deg(p^2)=750^2\cdot(-\frac{56}{5625})=-5600.
\]
It now suffices to compute $G_{\alpha_i\beta_i,h}(v_i,v_i)$, where these classes are given by the right-hand side of the decomposition \eqref{eq: referencing'}. We note that since $F$ has positive coefficients, we can readily verify the denominators appearing in $G_{\alpha_i\beta_i,h}(v_i,v_i)$ are positive. To simplify notations,  for a quadratic $Q$, we denote the Hessian matrix of $Q\cdot F$ by $H_Q$, then we have
\[
H_{X^2}=
\begin{pmatrix}
a&b&b\\
b&1&r\\
b&r&1
\end{pmatrix},\quad
H_{Y^2}=
\begin{pmatrix}
1&s&1\\
s&s^2&s\\
1&s&1
\end{pmatrix},\quad
H_{Z^2}=
\begin{pmatrix}
1&1&s\\
1&1&s\\
s&s&s^2
\end{pmatrix},
\]
and
\[
H_{XY}=
\begin{pmatrix}
b&1&r\\
1&s&1\\
r&1&1
\end{pmatrix},\quad
H_{XZ}=
\begin{pmatrix}
b&r&1\\
r&1&1\\
1&1&s
\end{pmatrix},\quad
H_{YZ}=
\begin{pmatrix}
r&1&1\\
1&s&1\\
1&1&s
\end{pmatrix}.
\]
Direct computation gives
\begin{gather*}
    H_{X^2}\begin{bmatrix}
        30\\-17\\-17
    \end{bmatrix}=\vec 0, \ \ H_{XY}\begin{bmatrix}
       10 \\ -2 \\ -9 
    \end{bmatrix}=\vec 0,\ \ 
    H_{XZ}\begin{bmatrix}
       10 \\ -9 \\ -2
    \end{bmatrix}=\vec 0,\\
    H_{Y^2}\begin{bmatrix}
       1 \\ 0 \\ -1
    \end{bmatrix}=H_{Y^2}\begin{bmatrix}
       0 \\ 2 \\ -1
    \end{bmatrix}=\vec 0,\ \ 
    H_{Z^2}\begin{bmatrix}
       1 \\ -1 \\ 0
    \end{bmatrix}=H_{Z^2}\begin{bmatrix}
       0 \\ -1 \\ 2
    \end{bmatrix}=\vec 0,
\end{gather*}

where the first equality implies $\deg(\chi^2(30\chi-17\gamma-17\zeta)w)=0$ for any $w\in \mathrm{span}\{\chi,\gamma,\zeta\}$, and in particular,
\[
    G_{\chi^2,h}(30\chi-17\gamma-17\zeta,30\chi-17\gamma-17\zeta)=0.
\]
The same argument works for other equalities, and we have that
\[
    \deg(p^2)+25\left(\sum_{i=1}^r\sqrt{G_{\alpha_i\beta_i,h}(v_i,v_i)}\right)^2=\deg(p^2)=-5600<0\leq \frac{2\deg(ph^2)^2}{\deg(h^4)}
\]
when $\epsilon=0$. Since the left-hand side is continuous in $\epsilon$, the denominators stay positive and the inequality stays true for $0<\epsilon\ll 1$.

\subsection{Verification of the Lorentzian property}
Now we show that $F$ is Lorentzian and that $F_\epsilon$ is strictly Lorentzian for $0<\epsilon\ll 1$. Since $F_{\epsilon}=F(Mx)$ where $M=I+\epsilon \begin{bsmallmatrix}
    1&1&1\\1&1&1\\1&1&1
\end{bsmallmatrix}$ is invertible with positive entries for $0<\epsilon \ll 1$, by \cref{prop: strictification} it suffices to prove that $F$ is concise and Lorentzian. Since $F$ has full support, for Lorentzian property, it remains to check that the various matrices $H_Q$ have at most one positive eigenvalue, according to \cref{thm: Lorentzian-criterion}.

For $Q=X^2,XY$ or $XZ$, we've seen that $H_Q$ has one null vector, given by the factors on the right-hand side of \eqref{eq: referencing'}, so it suffices to show that $H_Q$ is not positive semi-definite. This can be done by evaluating principal minors, and we have
\[
    \det(H_{X^2,\{2,3\}})=1-r^2<0,\ \ \det(H_{XY,\{2,3\}})=\det(H_{XZ,\{2,3\}})=s-1<0,
\]
where $H_{Q,I}$ denotes the principal submatrix of $H_Q$ by taking columns and rows in $I$. For $Q=Y^2$ or $Z^2$, there's nothing to check since the null space of $H_Q$ has dimension $2$. For $Q=YZ$, we have
\[
    H_{YZ}\begin{bmatrix}
        0\\1\\-1
    \end{bmatrix}=-\frac{1}{2}\begin{bmatrix}
        0\\1\\-1
    \end{bmatrix}
\]
and
\[
    \det(H_{YZ})=\frac{7}{40}>0,
\]
so $H_{YZ}$ has the Lorentzian signature.
    Now it remains to check conciseness. For any linear form $L$ such that $L\cdot F=0$, we must also have $L\cdot(YZ\cdot F)=0$. But we just found that $H_{YZ}$ is non-singular, therefore $L=0$.

\begin{remark}
    The first breakthrough of this project was a similar example given by the Danus system to illustrate that Lorentzian polynomials might fail to arise from the K\"ahler package. The original argument was to show that the element $p$ is truly primitive with respect to $h$ and therefore $\deg(p^2)$ should be nonnegative. The author later prompted the system to transform the argument into an inequality. The example presented here is also generated based on the same idea but is easier to present.
\end{remark}

\subsection{Comparison of types of polynomials}

We investigate different sets of polynomials. Let $L^d_n\subseteq S^d$ ($\mathring{L}^d_n$) denote the set of (strictly) Lorentzian polynomials in $n$ variables of degree $d$. According to \cref{thm: lorentzian and HR}, these are precisely the polynomials whose apolar algebras have mixed ($\text{HR}^{\leq 1}$) when nonzero. We follow the idea of \cite{MMS25} and make the following definition.

\begin{definition}
    A homogeneous polynomial $F\in S^d$ is said to be $k$-Lorentzian if $F=0$ or the triple $(A^\bullet_F,\deg,K_F)$ satisfies mixed ($\text{HR}^{\leq k}$). The set of all $k$-Lorentzian polynomials in $n$ variables in degree $d$ is denoted $L(k)^d_n$. Note that $L(1)^d_n=L^d_n$.
\end{definition}

\begin{definition}
    Let $\Bbbk$ be an algebraically closed field.
    \begin{enumerate}
        \item A homogeneous polynomial $F\in S^d$ of degree $d$ is called a \textit{realizable volume polynomial over $\Bbbk$} if there exists a $d$-dimensional projective variety $Y$ over $\Bbbk$ and a collection of semiample divisors $D_1,\ldots,D_n$ on $Y$ such that
    \[
        F=\frac{q}{d!}\int_Y\bigg(\sum_{i=1}^nx_iD_i\bigg)^d
    \]
    for some $q\in\Q_{\geq 0}$. The set of all such polynomials is denoted $\Bbb V^d_n(\Q,\Bbbk)$
        \item A homogeneous polynomial $F\in S^d$ is a \textit{volume polynomial over $\Bbbk$} if it is a limit of realizable volume polynomials over $\Bbbk$. The set of all such polynomials is denoted $\Bbb V^d_n(\R,\Bbbk)$.
    \end{enumerate}
\end{definition}

We introduce new families of polynomials as follows. Note that there's no need to invoke (mixed) hard Lefschetz, see for example \cite[Lemma 2.10]{MMS25}.

\begin{definition}
    \begin{enumerate}[leftmargin=*]
        \item A homogeneous polynomial $F\in S^d$ is a \textit{strict volume polynomial with respect to mixed $(\text{HR}^{\leq k})$} if there exists a Poincar\'e duality algebra $(A^\bullet,\deg,K)$ satisfying mixed $(\text{HR}^{\leq k})$, and there exists $\ell_1,\ldots,\ell_n\in K$, linearly independent, such that $F$ is the volume polynomial with respect to $\ell_1,\ldots,\ell_n$. The set of all such polynomials is denoted $\mathring{\Bbb V}^d_{\text{HR}^{\leq k},n}$
        \item A homogeneous polynomial $F\in S^d$ is a \textit{volume polynomial with respect to mixed $(\text{HR}^{\leq k})$} if there exists a Poincar\'e duality algebra $(A^\bullet,\deg,K)$ satisfying mixed $(\text{HR}^{\leq k})$, and there exists $\ell_1,\ldots,\ell_n\in \overline K$ such that $F$ is the volume polynomial with respect to $\ell_1,\ldots,\ell_n$. The set of all such polynomials is denoted $\Bbb V^d_{\text{HR}^{\leq k},n}$. We also let $\Bbb {V'}^{d}_{\text{HR}^{\leq k},n}$ denote the subset of $\Bbb {V}^{d}_{\text{HR}^{\leq k},n}$ such that the algebra can be chosen so that $n\leq\dim A^1$. Note that this automatically includes all concise polynomials in $\Bbb {V}^{d}_{\text{HR}^{\leq k},n}$.
    \end{enumerate}
\end{definition}

\begin{proposition}\label{prop: volume and lorentzian}
    A (strict) volume polynomial with respect to mixed $(\text{HR}^{\leq 1})$ is the same as a (strictly) Lorentzian polynomial.
\end{proposition}

\begin{proof}
    One direction is supplied by \cref{prop: strictification}.

    For the other direction, \Cref{thm: lorentzian and HR} shows that $\Bbb V^d_{\text{HR}^{\leq 1},n}\supseteq L^d_n$. Given a strictly Lorentzian polynomial $F(x)$, the polynomial $G(x)=F((I+\epsilon E)^{-1}x)$ is also strictly Lorentzian for $0<\epsilon\ll1$, where $E$ is the all-ones matrix. It's clear that in the apolar algebra of $G$, the volume polynomial with respect to $X_1+\epsilon(X_1+\cdots+X_n),X_2+\epsilon(X_1+\cdots+X_n),\ldots,X_n+\epsilon(X_1+\cdots+X_n)$ is exactly $F$, which shows that $F$ is a strict volume polynomial with respect to mixed $(\text{HR}^{\leq 1})$.
\end{proof}

We have the following comparison of different notions.
\begin{proposition}
    We have the following inclusions for any $1\leq k\leq \lfloor\frac{d}{2}\rfloor$.
    \[\begin{tikzcd}
	 & & L^d_n & \\
	 \Bbb V^d_n(\R,\Bbb C) & \overline{\mathring{\Bbb V}^d_{\text{HR}^{\leq k},n}} & {\Bbb V}^d_{\text{HR}^{\leq k},n} & L(k)^d_n \\
    \Bbb V^d_n(\Q,\Bbb C) & {\Bbb V'}^d_{\text{HR}^{\leq k},n} & & \\
	& \mathring{\Bbb V}^d_{\text{HR}^{\leq k},n}
	\arrow[hook, from=2-1, to=2-2]
	\arrow[hook, from=2-2, to=1-3]
	\arrow[hook, from=2-3, to=1-3]
	\arrow[hook', from=2-4, to=2-3]
	\arrow[hook, from=3-1, to=2-1]
    \arrow[hook, from=3-1, to=3-2]
	\arrow[hook, from=3-2, to=2-2]
	\arrow[hook, from=3-2, to=2-3]
	\arrow[hook, from=4-2, to=3-2]
\end{tikzcd}\]
All inclusions here, except for the ones starting from ${\Bbb V'}^d_{\text{HR}^{\leq k},n}$, are strict in general.
\end{proposition}

\begin{proof}
    To check the inclusion from the left column to the middle column, given a realizable volume polynomial, after blowing up the variety $Y$ enough times, one can take $\bigoplus_kH^{k,k}(Y)\cap H^\bullet(Y,\R)$ which has the K\"ahler package and a sufficiently large degree $1$ component. Other inclusions directly follow from the definition. The vertical inclusion on the left is strict simply because a realizable volume polynomial cannot have irrational coefficients.

    According to \cite[Footnote 15]{BH20}, there exists a concise cubic Lorentzian polynomial that fails to satisfy the \textit{reverse Khovanskii-Teissier inequality}, hence fails to be a volume polynomial over any field. On the other hand, \cref{prop: volume and lorentzian} shows that it is a volume polynomial with respect to mixed $(\text{HR}^{\leq 1})$, so the inclusions from the left column to the middle column are strict when $k=1$.
    
    The existence of concise non-strict Lorentzian polynomial, for example $xy$, shows the strictness of the vertical inclusion at the bottom when $k=1$.

    As we've seen in \cref{ex: ternary}, a strictly Lorentzian polynomial (and its small perturbations) might not be realized as the volume polynomial in any ring with mixed $(\text{HR}^{\leq 2})$, which shows the strictness of the inclusions $\overline{\mathring{\Bbb V}^d_{\text{HR}^{\leq k},n}} \subseteq L^d_n$ and ${\Bbb V}^d_{\text{HR}^{\leq k},n} \subseteq L^d_n$ when $k\geq 2$.

    \cite[Example 6.17]{NW26} shows that there exist binary quartic realizable volume polynomials whose apolar algebra doesn't satisfy mixed $(\text{HR}^{\leq 2})$, which means the inclusions $L(k)^d_n\subseteq {\Bbb V}^d_{\text{HR}^{\leq k},n}$ and $L(k)^d_n\subseteq L^d_n$ are strict for $k\geq 2$.
\end{proof}

\begin{remark}
\begin{enumerate}[leftmargin=*]
    \item The introduction of ${\Bbb V'}^d_{\text{HR}^{\leq k},n}$ concerns whether one can enlarge an algebra with the K\"ahler package without changing its degree. When $k=\lfloor\frac{d}{2}\rfloor$, one can consider the algebras
    \[
        B\coloneqq A\otimes \R[x]/x^2,\quad C\coloneqq B/\ann(h+x),
    \]
    where $h\in K$. By \cite[Proposition 7.7]{AHK18} and \cite{Cat08}, $B$ has the K\"ahler package with respect to $(K+\R_{>0}x)$. Therefore, by the descent lemma \cref{prop: descent}, $C$ also has the K\"ahler package with respect to the image of $(K+\R_{>0}x)$. When $d\geq 3$, we have
    \[
        \dim C^1=\dim A^1+1,\qquad \deg_C(\bar \alpha_1\cdots\bar \alpha_d)=\deg_A(\alpha_1\cdots\alpha_d),
    \]
    for elements $\alpha_1,\ldots\alpha_d\in A^1$. This shows that ${\Bbb V}^d_{\text{HR}^{\leq k},n}={\Bbb V'}^d_{\text{HR}^{\leq k},n}$ and we can compare $\overline{\mathring{\Bbb V}^d_{\text{HR}^{\leq k},n}}$ and ${\Bbb V}^d_{\text{HR}^{\leq k},n}$ directly. For general $k$, it remains open whether this ``blowup'' construction still works.
    \item When $k=1$, \Cref{prop: volume and lorentzian} shows that $\overline{\mathring{\Bbb V}^d_{\text{HR}^{\leq k},n}}$ and ${\Bbb V}^d_{\text{HR}^{\leq k},n}$ are equal. Whether this remains true when $k>1$ remains open.
\end{enumerate}
\end{remark}